\chapter{Os Cinco Khandhas}

\lettrine{E}{nquanto estes corpos humanos} estão vivos e seus sentidos físicos estão a
funcionar, temos que estar constantemente em guarda, alertas e conscientes,
porque a força do hábito de apreender o mundo sensual como um eu é muito forte.
Este é um condicionamento muito forte em todos nós. Portanto, o caminho que o
Buddha ensinou é o caminho da consciência e da sábia reflexão. Em vez de fazer
declarações metafísicas sobre as nossas Verdadeiras Naturezas ou Realidade
Última, os ensinamentos do Buddha apontam para a condição do apego (a posse).
Essa é a única coisa que nos impede de alcançar a iluminação.

A sabedoria de Buddha é a compreensão de como as coisas são pela via da
observação de nós mesmos, em vez de observarmos apenas como as estrelas e os
planetas funcionam. Não saímos olhando para as árvores e contemplando a Natureza
como se ela fosse um objecto da nossa visão, mas na realidade estamos a observar
a Natureza como ela funciona através desta formação pessoal.

O que consideramos ser pode ser classificado como cinco agregados ou
\emph{khandhas}: \emph{rūpa}, forma; \emph{vedanā}, sensação/sentimento;
\emph{saññā}, percepção; \emph{saṅkhārā}, formação mental ou processo de pensamento, e
\emph{viññāṇa}, consciência dos sentidos (cognição). Estes providenciam um meio
apto para ver todos os fenómenos sensuais em grupos. O mais fácil para se
meditar é o \emph{khandha rūpa}, a forma do próprio corpo, porque podemo-nos
sentar aqui -- ele está assente no chão, pesado, grosseiro. É uma coisa em um
movimento mais lento do que os fenómenos mentais -- \emph{vedanā}, \emph{saññā},
\emph{saṅkhārā} ou \emph{viññāṇa}. Podemos realmente reflectir sobre o nosso
próprio corpo por longos períodos de tempo, meditar sobre a respiração em vez de
sobre a consciência, porque a respiração é algo em que nos podemos concentrar.
% REVISÃO: "em um movimento mais lento" (acima) — em PT-PT contrai-se: "num movimento mais lento".
% REVISÃO: "contemplar sua própria respiração" — falta o artigo: "contemplar a sua própria respiração".
Pessoas comuns podem contemplar sua própria respiração.

Podemos contemplar a sensação dos nossos próprios olhos. Eles possuem sensações.
Contempla-se a língua, a humidade da boca ou a língua tocando o céu da boca.
Podemos contemplar o corpo como um órgão dos sentidos, dando-nos sensações de
prazer e de dor, calor e frio. Observe-se simplesmente como é a sensação de frio
ou de calor no corpo; podemos contemplar isso porque não é o que nós
\emph{somos}. É um objecto que facilmente se pode ver e observar como se fosse
algo separado de nós. Se não fizermos isso, temos tendência para apenas reagir.
Quando estamos com muito calor, tentamos arrefecer e tirar a camisola. Então
ficamos com frio e vestimo-la novamente. Podemos simplesmente reagir a essas
sensações de prazer e dor no corpo. Prazer: ``Oh, que maravilhoso'', tentamo-nos
agarrar a isso, ter mais prazer. E a dor: ``Oh'' -- livre-se disso; foge-se de
qualquer coisa desagradável ou dolorosa. Mas na meditação nós podemos ver estas
sensações e que o próprio corpo é uma condição sensual que possui prazer, dor,
calor e frio.

Podemos reflectir sobre as formas que vemos. Basta olhar para algo bonito, como
as flores. As flores são provavelmente as coisas mais bonitas na terra, por isso
gostamos de flores. Portanto, observemos quando olhamos para uma flor, como
somos atraídos por ela, e como queremos continuar a olhar para ela: ser atraído
pelo que é agradável aos olhos. Ou, então, algo desagradável aos olhos. Olhar
para, por exemplo, excrementos. Quando vemos excrementos, esterco de vaca no
caminho, ignoramos educadamente. Olhemos para o nosso próprio excremento. Nós
mesmos o produzimos e, no entanto, é algo que na realidade não queremos andar
por aí a mostrar às outras pessoas. Preferíamos que nunca ninguém nos visse a
fazê-lo. Na verdade, não nos sentimos atraídos a olhar para isso como para uma
flor, pois não? E no entanto, estamos dispostos a usar flores, a transportá-las
e a mantê-las no nosso altar.

Não se trata de que devíamos achar o excremento atraente. Estou apenas a apontar
para o facto de que podemos meditar sobre esta força do mundo sensorial. É uma
força natural. Não é mau ou errado, mas podemos meditar sobre isso a fim de
vermos a nossa tendência de reagir às experiências sensoriais.

Quando ouvimos sons bonitos ou sons horríveis, odores agradáveis ou fedorentos,
sabores agradáveis ou desagradáveis, sensações físicas prazerosas ou dolorosas,
calor ou frio -- meditemos sobre elas. Olhemos e vejamos estas coisas como elas
são: todo o \emph{rūpa} é impermanente. Flores bonitas são só bonitas por um
tempo; a partir de determinado momento tornam-se repulsivas. Portanto, estamos a
observar esta transformação natural do que é novo e bonito para o que é velho e
feio. Eu mesmo era muito mais bonito aos vinte anos. Agora estou velho e feio.
Um corpo humano velho não é muito bonito, certo? Mas é o corpo, a percorrer o
que é suposto. Estou feliz que não esteja a ficar mais bonito. Seria embaraçoso
se fosse.

Ora, os \emph{khandhas} mentais também operam com base no mesmo princípio.
\emph{Vedanā} é um estado mental, o sentimento que se tem da atração e aversão
em relação às coisas físicas que se ouve, vê, cheira, prova ou toca. A sensação
de dor é exactamente como é, mas depois há a reacção de gostar ou desgostar. Ou
até nem isso, mas apenas um movimento na sua direcção ou para longe dessa
sensação.

Podemos estar cientes das sensações, dos humores. Observe-se o calor que vem da
raiva, a apatia que vem da dúvida e da preguiça e torpor. Observe-se a
sensação de quando estamos com ciúmes. Podemos testemunhar essa sensação.
Observemos, em vez de apenas tentar aniquilar o ciúme. Quando o ciúme condiciona
a nossa mente, em vez de reagirmos ou tentarmos livrarmo-nos do ciúme por não
gostarmos dele, podemos começar a reflectir sobre ele. Quando estamos com frio,
o que é a natureza do frio? Gostamos disso? Esta frieza, sentir frio, é algo
% REVISÃO: a frase é interrogativa ("...ou será que beneficiamos muito disso?") — falta o ponto de interrogação.
terrivelmente desagradável ou será que beneficiamos muito disso. Fome, como é a
fome? Quando sentirmos fome, meditemos na sensação física à qual temos tendência
para reagir tentando obter algo para comer.

Ou meditemos sobre a sensação de se estar sozinho ou separado, ou a sensação de
que as pessoas nos desprezam. Se acharem que eu não gosto de vocês -- meditem
sobre esse sentimento. Ou se não gostam de mim! Meditem sobre isso. Tragam isso
para a consciência agora. Não analiticamente, tentando perceber o que realmente
sinto; ou se o vosso relacionamento comigo é um relacionamento de dependência
infantil que não deveriam ter; ou envolvendo-se na psicologia freudiana ou o que
quer que seja. Mas apenas observem o estado de espírito de dúvida e incerteza
nos vossos relacionamentos com os outros -- não para analisar, mas
apenas para observar o tipo de sentimentos de confiança ou falta de confiança,
aversão ou atracção. Isso é \emph{vedanā}. Isto é uma coisa natural. Somos todos
seres sensíveis, então há uma atracção e repulsão a ocorrer o tempo inteiro. É
uma condição da natureza, não um problema pessoal -- a menos que o tornemos
pessoal.

\emph{Saññā khandha} é o \emph{khandha} percepção. Compreender uma percepção
significa acreditar na maneira como as coisas aparecem no presente como se fosse
uma espécie de qualidade permanente. É assim que temos tendência para funcionar
nas nossas vidas. Então, eu posso pensar, por exemplo, ``O venerável Viradhammo
é assim''. É uma percepção que tenho quer eu esteja aqui, sentado ao lado do
Venerável Viradhammo, ou sozinho, quer ele esteja a ajudar-me ou com raiva de
mim. Eu tenho esta forma de ver fixa. Uma percepção fixa não é muito consciente,
mas se eu acreditar na minha percepção tenho tendência para funcionar dessa
posição fixa específica. % REVISÃO: "se sua personalidade" — falta o artigo (PT-PT): "se a sua personalidade".
Dessa forma, quando penso nele, é como se sua
personalidade fosse fixa e constante, em vez de ser do modo que é agora. A minha
percepção dele é apenas uma percepção do momento; não é uma alma que atravessa o
tempo nem é uma personalidade fixa. Portanto, deve-se meditar sobre
\emph{saññā}.

% REVISÃO: quebra de parágrafo indevida entre "quer" e "nos consideremos" — remover a linha em branco; é uma só frase.
\emph{Saṅkhārā} são formações mentais. Existem percepções da mente
(\emph{saññā}) e funcionamos a partir delas (\emph{saṅkhārā}). Então, as
presunções que temos acerca de nós mesmos -- da infância, pais, professores,
amigos, parentes e tudo mais; quer

nos consideremos bons e positivos ou de uma forma negativa ou confusa -- isso é
tudo os \emph{khandhas} \emph{saññā}/\emph{saṅkhārā}.

As memórias surgem, ou o tipo de medos sobre o que nos poderá estar a faltar.
Podemos preocupar-nos com a possibilidade de haver uma falha séria no nosso
carácter ou certos desejos horríveis reprimidos à espreita no fundo da nossa
mente -- que podem surgir na meditação e deixar-nos loucos! \emph{Isso} é outra
condição mental, esse não saber o que somos. Então por vezes imaginamos as
piores coisas possíveis. Mas o que podemos saber é que tudo o que
\emph{acreditamos}, é uma condição da mente: surge, desaparece, é impermanente!

Se agirmos sob certas percepções fixas que temos de nós mesmos, então concebemos
todos os tipos de coisas. Se agimos a partir da perspectiva ``Eu sou um homem''
e logo nos tornamos essa percepção de nós mesmos, assumimos isso. Então, nunca
investigaremos essa percepção; apenas acreditamos que ``sou um homem'' e então
concebemos masculinidade como sendo de certa maneira; ``o que um homem deveria
ser''. Depois comparamo-nos com aquilo que é o ideal para a masculinidade e
quando não vivemos à altura desses altos padrões de masculinidade,
preocupamo-nos. Algo errado! Começamos a sentirmo-nos chateados, culpados ou a
odiarmo-nos por causa da presunção básica de ser um homem. A um nível
convencional, isto pode ser verdade; os homens são desta maneira e as mulheres
daquela. Não estamos a negar a realidade convencional, mas já não nos estamos a
apegar a ela como uma qualidade pessoal, uma posição fixa para assumirmos em
todos os momentos e lugares.

Esta é uma forma de nos libertarmos da qualidade que se vincula às condições
insatisfatórias. Porque se acreditamos que somos homem ou mulher como nossa
verdadeira identidade e nossa alma, isso irá então conduzir-nos sempre a um
estado de depressão mental. Tudo isso são percepções que temos. Criamos muita
miséria por nos vermos como pretos ou brancos, ou membros de uma determinada
nacionalidade ou classe. Na Inglaterra, as pessoas sofrem com essa percepção de
pertencer a uma determinada classe; na América, sofremos por não termos
quaisquer percepções de classe, com a percepção de que somos todos a mesma
coisa, somos todos iguais. É o apego a qualquer uma destas, mesmo à percepção
mais elevada e mais igualitária, que nos conduz ao desespero.

Investigando esses cinco ``aglomerados''\footnote{O significado literal de
  \emph{khandhas}}, agregados ou grupos, começamos a percebê-los. Podemos
conhecê-los como objectos porque são \emph{anatta} -- não-eu. Se eles fossem o
% REVISÃO: "seriamos" — falta o acento: "seríamos".
% REVISÃO: gentílicos ("Italiano, Dinamarquês, Suíço, Inglês ou Americano", abaixo, e "Americano" mais adiante) escrevem-se com minúscula em português: italiano, dinamarquês, suíço, inglês, americano.
que somos, não conseguiríamos vê-los. Só seriamos capazes de \emph{ser} eles.
Não teríamos forma de testemunhá-los ou de nos desprendermos deles, apenas
ficaríamos presos dentro deles o tempo inteiro, sem qualquer habilidade de nos
desprendermos e observá-los. Mas sermos homens, mulheres, monges, freiras,
Italiano, Dinamarquês, Suíço, Inglês ou Americano, ou qualquer outra coisa, é
apenas uma verdade relativa, relativa a certas situações.

Ainda assim, agimos nas nossas vidas a partir de posições fixas, de ser ``Eu sou
Americano'' e ``Nós somos assim''. Em todo o mundo, temos preconceitos nacionais
e raciais. Estes são apenas percepção e concepção (\emph{khandhas
  saññā/saṅkhārā}) que podemos observar. Quando possuímos uma visão fixa sobre
alguém -- ``Uma coisa que eu não suporto são os Hondurenhos'' -- podemos
observar isso na nossa mente, certo? Mesmo que tenhamos fortes preconceitos e
sentimentos e tentemos livrar-nos deles, isso resulta da presunção de que
\emph{não devemos} ter quaisquer preconceitos e não devemos ter quaisquer maus
sentimentos em relação a ninguém e que \emph{devíamos} ser capazes de aceitar
críticas com uma mente equânime e não ficarmos com raiva ou chateados. Essa é
outra presunção muito idealista, não é? Percebemos isto como uma condição da
mente e continuamos a observar. Em vez de nos odiarmos ou odiarmos os outros por
serem preconceituosos, observamos precisamente as limitações de quaisquer
preconceitos ou percepções e concepções da mente. Meditamos sobre a natureza
impermanente da percepção. Por outras palavras, não estamos a tentar justificar
ou explicar ou vermo-nos livres de ou mudar nada. Estamos simplesmente a tentar
observar que todas as coisas mudam -- tudo o que começa acaba.

Em seguida, meditamos no \emph{khandha viññāṇa}, a consciência (cognitiva)
sensorial, consciência sensorial do olho, ouvido, nariz, língua, corpo e mente.
Como uma coisa vai para a outra, ciente dos movimentos da consciência dos
sentidos. Olhando para algo, ouvindo algo -- muda muito rapidamente.

Todos estes cinco \emph{khandhas} são \emph{aniccā}, impermanentes. Quando
cantamos: \emph{rūpaṁ aniccaṁ, vedanā aniccā, saññā aniccā, saṅkhārā aniccā,
  viññāṇaṁ aniccaṁ}, isto é muito profundo. Depois: \emph{sabbe saṅkhārā aniccā}.
  \emph{Saṅkhārā} significa ``todos os fenómenos condicionados'', todas as
experiências sensoriais -- os órgãos dos sentidos, os objectos dos órgãos dos
sentidos, a consciência que surge no contacto -- tudo isto é \emph{saṅkhārā} e é
\emph{aniccā}. Tudo está condicionado. Portanto, \emph{saṅkhārā} inclui os
outros quatro: \emph{rūpa, vedanā, saññā, viññāṇa}.

Com isto, temos uma perspectiva a partir da qual o mundo condicionado é
infinitamente variável e complexo. Mas onde é que separamos \emph{saññā} de
\emph{saṅkhārā} e \emph{saṅkhārā} de \emph{viññāṇa} e tudo isso? É melhor não
tentar obter divisões precisas entre estes cinco aglomerados; eles são apenas
formas convenientes de se ver as coisas, ajudando-nos a meditar sobre os estados
mentais, o mundo físico e o mundo sensorial.

Não estamos a tentar consertar nada como ``Isto é permanentemente
\emph{saṅkhārā} e aquilo é definitivamente \emph{saññā}'', mas sim a usar essa
discriminação para observar que o mundo sensorial -- do físico ao mental, do
grosseiro ao subtil -- é condicionado e que todos os fenómenos condicionados são
impermanentes. Então temos uma maneira de ver a totalidade do mundo condicionado
como impermanente, em vez de nos envolvermos em tudo. Nesta prática da meditação
introspectiva, não estamos a tentar analisar o mundo condicionado, mas a
destacarmo-nos dele, para vê-lo em perspectiva. Isto é quando realmente
começamos a compreender \emph{anicca}; sabemos introspectivamente \emph{sabbe
  saṅkhārā aniccā}.

Portanto, quaisquer pensamentos e crenças que temos são apenas condições. Mas
não estou a dizer que não devemos acreditar em nada, estou apenas a apontar uma
maneira de ver as coisas em perspectiva, para evitar sermos iludidos por elas.
Não obteremos a experiência do vazio ou do Incondicionado, a Realidade Imortal,
como uma realização pessoal. Alguns de vós têm tentado realizar isso como um
tipo de realização pessoal, certo? ``Eu conheço o vazio. Eu percebi o vazio'' --
dando pancadinhas nas próprias costas. Isso não é \emph{sabbe dhammā anattā} --
isso é querer agarrar o Incondicionado, torná-lo uma condição. ``Eu'' e ``Meu''.
Quando começamos a pensar em nós próprios como tendo percebido o vazio, podemos
ver isso igualmente como uma condição da mente.

Agora \emph{sabbe dhammā anattā}: todas as coisas são não-eu, não uma pessoa,
não uma alma permanente, não é um eu de qualquer género. Isso também é muito
importante contemplar porque \emph{sabbe dhammā} inclui todas as coisas; os
fenómenos condicionados do mundo sensorial e o Incondicionado, a Realidade
Imortal.

Note-se igualmente que os budistas não mencionam a Realidade Imortal como sendo
um eu! ``Eu tenho uma alma imortal'' ou ``Deus é a minha verdadeira natureza!''
O Buddha evitou quaisquer declarações desta natureza. Qualquer tentativa de nos
concebermos como qualquer coisa que seja, é um obstáculo à iluminação porque
ainda nos agarramos a uma ideia, a um conceito do eu como parte de algo. Talvez
pensemos que há um pedaço de nós, uma pequena alma que se junta àquela maior
quando se morre. Essa é uma concepção da mente -- não é? -- que se pode
conhecer. Não estamos a dizer que é mentira ou falso, mas estamos apenas a ser o
saber, saber o que pode ser conhecido. Não nos sentimos compelidos a apreender
isso como uma crença; vemos isso apenas como algo que vem da mente, uma condição
da mente, então até isso deixamos ir.

Mantenhamos essa fórmula ``todas as condições são impermanentes, todas as coisas
são não-eu'' para reflexão. E então na vida, à medida que a vivemos, aconteça o
que acontecer, conseguimos perceber \emph{sabbe saṅkhārā aniccā, sabbe dhammā
  anattā}. Isso evita que nos iludamos com fenómenos milagrosos que nos podem
acontecer e é uma forma de compreendermos outras convenções religiosas. Os
cristãos vêm e dizem: ``Só por meio de Jesus Cristo se pode ser salvo. Você não
pode ser salvo pelo Budismo. Buddha era apenas um homem, mas Jesus Cristo era o
Filho de Deus.'' Então você pensa ``Oh, questiono-me, talvez eles estejam
certos.'' Afinal, quando se vai a um desses encontros de ``Nascidos de Novo'',
todos irradiam felicidade; os seus olhos estão brilhantes e dizem: ``Louvado
seja o Senhor!'' Mas quando se vai a um mosteiro budista e ficamos ali
simplesmente sentados, horas a fio, observando a respiração, não ficamos logo em
êxtase. Começamos a duvidar e pensamos: ``Talvez isso esteja certo; talvez Jesus
seja o caminho''. Mas o que \emph{podemos} saber é que há uma dúvida. Observemos
essa dúvida ou o sentimento de sermos intimidados por outras religiões quando
elas vêm em força, ou o sentimento de aversão relativamente a elas ou ter
preconceitos contra religiões. O que podemos saber é que essas são percepções da
mente: elas vêm, vão e mudam.

Mantenhamos uma reflexão constante e calma sobre estas coisas, em vez de
tentarmos descobrir ou sentir que precisamos de justificar o facto de sermos
budistas. Os cristãos começam a dizer: ``Vocês não fazem nada pelo terceiro
mundo'', e dizemos, ``Nós\ldots{} nós\ldots{} nós\ldots{} recitamos os cânticos!
Partilhamos mérito e irradiamos bondade amorosa.'' Isto parece fraco numa
situação em que falamos de desnutrição e fome em África! Mas agora, neste
momento, existe a oportunidade de entendermos os limites para o que podemos
fazer. Todos nós, se pudéssemos, definitivamente faríamos algo relativamente à
fome em África se soubéssemos que existia algo que pudéssemos fazer aqui e
agora, neste momento. Reflictam sobre isso -- qual é o problema na realidade? É
a fome em África ou é o egoísmo e a ignorância humanos? A fome em África não é o
resultado da ganância humana, egoísmo e estupidez?

Portanto, abrimos as nossas mentes ao Dhamma. Reflectimos sabiamente
sobre isso e então percebemos. A verdade deve ser percebida e conhecida dentro
do contexto da experiência pessoal\ldots{} Mas a prática é contínua -- após 18
anos, ainda pratico o tempo inteiro. As coisas mudam: as pessoas elogiam e
culpam e o mundo continua. Continuamos simplesmente a reflectir sobre isso com
\emph{sabbe saṅkhārā aniccā, sabbe dhammā anattā}. Quando reconhecemos o
condicionado e o Incondicionado, então temos o que se chama a capacidade para
desenvolver o Caminho e não há mais confusão sobre isso.

Agora o objectivo é realizar o Nibbāna, ou Realidade Imortal, ou o
desapego -- perceber como é não estarmos apegados aos cinco \emph{khandhas}.
Percebamos isso quando estivermos aqui sentados e realmente em paz. Então não
haverá apego aos cinco \emph{khandhas}; mas pode acontecer gerarmos uma
percepção dessa paz plena e apegarmo-nos a isso e tentarmos sempre meditar para
ficarmos em paz de novo de acordo com essa percepção. É por isso que a prática é
mais um contínuo abrir mão do que uma conquista.

Às vezes, quando ficamos calmos nestes retiros, temos uma mente muito tranquila.
E apegamo-nos; então meditamos para alcançar esse estado de bem-aventurança. Mas
a meditação da compreensão introspectiva significa investigarmos a natureza das
coisas, dos cinco \emph{khandhas}, vendo-os como \emph{anicca} -- impermanentes;
como \emph{dukkha} -- insatisfatórios. Nenhum destes \emph{khandhas} tem a
capacidade de nos dar qualquer tipo de satisfação permanente. A sua própria
natureza é insatisfatória e \emph{anatta}.

Comecem a investigar e sabiamente considerem \emph{sabbe saṅkhārā aniccā, sabbe
  dhammā anattā}, em vez de pensarem que atingiram algo ou que têm que se
segurar a essa realização, e começarem a ressentir-se de qualquer pessoa que se
intrometa no vosso caminho. \emph{Percebam} o que é apego. Quando a vossa mente
estiver realmente concentrada, larguem isso. Em vez de simplesmente se
gratificarem nessa sensação de paz, agarrem-se a algo. Preocupem-se com algo.
Façam isso deliberadamente para começarem a ver como saem e se agarram a coisas
ou se preocupam em perder essa sensação de paz.

À medida que começamos a compreender e experienciar o abrir mão na nossa
prática, começamos a compreender o que os Buddhas sabem: \emph{sabbe saṅkhārā aniccā,
  sabbe dhammā anattā}. Não são apenas as palavras -- até um papagaio pode dizer
as palavras -- mas o papagaio não é iluminado, será? A compreensão interior é
diferente do conhecimento conceitual. Mas agora estamos a penetrar, indo a fundo
nisto, rompendo a ilusão do eu como sendo o que seja; ou nada -- se acreditamos
que não temos um eu -- essa é outra crença; ``Eu acredito que não tenho um eu''.

O Buddha apontou para o caminho entre estes dois extremos: acreditar que temos
um eu e acreditar que não temos um eu. Não se consegue encontrar nada nos cinco
\emph{khandhas} que tenha um eu ou alma permanente: as coisas surgem do
Incondicionado e retornam ao Incondicionado. Assim, é mais por abrir mão do que
por adoptar qualquer outra atitude, que não procuramos agarrar-nos mais às
condições mortais.

\photoPage{image-004-robe-cloth-measuring.jpg}{Bhikkhus a costurar os seus hábitos}{}
